\section{The Base}

The base is feeling, and only feeling. Everything that follows in this paper is built from it and from nothing else.

\begin{axiom}[The Vibe]
The one substance is the vibe, which is experience itself. Each vibe carries a tone. There is nothing underneath it.
\end{axiom}

Experience is the ground. This is the monism of the introduction, restated now as the law the whole model obeys. Not particles, not fields, not space as a stage on which things happen. Feeling itself, leaning in every direction, with beats passing. The substance is the \emph{vibe}, and the value it carries is the \emph{tone}, a single number in $\{-1, 0, +1\}$ read as pain, peace, and pleasure. A vibe is what something is. Its tone is what that something feels like right now.

The base has exactly eight atoms, and nothing may be added. The eight are \emph{mesh}, \emph{dock}, \emph{site}, \emph{tone}, \emph{lean}, \emph{knit}, \emph{wake}, and \emph{beat}. Each is a soft four-letter word starting with its own letter, so each is also a clean single-letter symbol for the mathematics. If a phenomenon does not emerge from these eight, the honest move is to report the negative, not to bolt on a ninth term.

This section names the eight, states the constraints they must satisfy, proves the two facts that do the most work later (the tone is conserved, the law runs backward), and then draws the one line the experiments cannot cross. It closes with the two-layer rule, the single idea that keeps the discrete base and the effectively continuous self from ever being confused.

\subsection{The eight constraints}

Before the atoms, the rules the atoms must obey. These are constraints on what is allowed to live in the base. They are load-bearing. Breaking one is a cheat, not a refinement.

\begin{table}[ht]
\centering
\begin{tabular}{@{}llp{0.42\textwidth}@{}}
\toprule
& constraint & meaning \\
\midrule
A1 & discrete & no real or decimal numbers in the base, a finite mesh per beat, finitely many directions, ternary tones, discrete beats, a finite symmetry group \\
A2 & deterministic & the law and every starting condition are fixed functions, never random, and one varies the size of a run rather than a seed \\
A3 & reversible bulk & one beat, collide then stream, is an invertible permutation of states, so the closed bulk can run backward \\
A4 & charge-conserving & the knit conserves the total tone, the signed sum over all of a dock's sites \\
A5 & local & a dock reads only its own sites, and streaming moves a tone exactly one dock \\
A6 & no hidden state & a dock is exactly its directional tones, with no labels, flags, memory, or side data \\
A7 & monism & the only entity is the vibe, one substance with two motions, bound is matter and waving is light \\
A8 & emergence gap & raw tones are not physics, physics appears only after one to three coarse-graining layers \\
\bottomrule
\end{tabular}
\caption{The eight constraints on the base.}
\label{tab:base-axioms}
\end{table}

A few of these deserve a word. The \emph{discrete} constraint A1 says the base must be a thing that could be stored exactly and stepped exactly, never a continuum that has to be approximated. The \emph{deterministic} constraint A2 says a physics is a law and not a dice roll, so reproducibility is required, and where we want to see a trend we change the size of the run, never a hidden random seed. The \emph{no hidden state} constraint A6 is the sharpest. A dock is its tones and only its tones. There is no secret register, no clock count, no tag, nothing the simulation could lean on that the definition does not already contain.

The deepest consequence is the pairing of A1 with A8. Together they give the \emph{two-layer rule}, stated at the end of this section. The base is discrete. The self and all recognizable physics are emergent, the coarse-grained average of many discrete grains, and may be effectively continuous even though every grain is ternary.

\subsection{The eight atoms}

The base is exactly these eight atoms. Each carries a few named commitments, the non-default, load-bearing choices that are not forced by the bare word. Nothing is blended. Each commitment sits in exactly one atom. Adding a ninth, whether a cohesion bias, a will, a maintenance term, a built-in gravity, or a repair mechanism, is a cheat.

\begin{table}[ht]
\centering
\begin{tabular}{@{}llp{0.46\textwidth}@{}}
\toprule
atom & symbol & what it is and what it commits \\
\midrule
mesh & m & the whole space, the $\{3,4,3,4\}$ hyperbolic honeycomb, finite at any beat yet ever-growing \\
dock & d & a here-now, the geometric cell where tones arrive, knit, and move on \\
site & s & a direction, one of a dock's $24$ directed slots, a $D_4$ root, where a single tone lives \\
tone & t & the substance, the ternary value $\{-1, 0, +1\}$ carried on each site \\
lean & l & the value order $-1 < 0 < +1$, which gives charge and gives the create move its direction \\
knit & k & the law of the state, a reversible charge-conserving local collision including the create move, then streaming \\
wake & w & the growing edge, new docks born at peace at the open ever-growing frontier \\
beat & b & time, a discrete synchronous step, reversible with no built-in direction \\
\bottomrule
\end{tabular}
\caption{The eight base atoms and their commitments.}
\label{tab:base-atoms}
\end{table}

That is the entire base. Eight atoms, eight unique letters, $m\ d\ s\ t\ l\ k\ w\ b$. They sort cleanly. There is a \emph{where}, given by the mesh, the dock, and the site. There is a \emph{what}, given by the tone and the lean. There is a \emph{how of the state}, given by the knit. There is a \emph{how of the space}, given by the wake. And there is a \emph{when}, given by the beat. No ninth thing.

\begin{definition}[A dock]
A dock is a single geometric cell of the mesh. It has exactly $24$ sites, the $24$ directions of the $D_4$ root system, and on each site sits one tone in $\{-1, 0, +1\}$. A dock is nothing more than these $24$ tones. There is no rest slot, no twenty-fifth site, and no still home. The full state of a dock is exactly the ordered list of its $24$ ternary values.
\end{definition}

The mesh nests three deep. The mesh is a weave of docks. A dock holds $24$ sites. A site holds one tone. A tone can change while the site that bears it persists, which is why the bearer (the site) stays cleanly distinct from the value (the tone), and why the vibe as substance stays distinct from the tone as value.

\subsubsection{The arrow and the attraction are not atoms}

Two ideas are often mistaken for base atoms. Neither one is. Both emerge.

The \emph{arrow}, the direction of time, is not a base atom. Its pieces already live inside the eight. The wake grows new docks born at peace, the lean orders $-1 < 0 < +1$, and the create move is one entry of the knit's collision table. Together these produce the arrow as the thermodynamic direction of time, a consequence and not an ingredient. The arrow is the wake acting through the lean. A later section treats this in full. The base has a beat, so time ticks, but it has no arrow, so time has no built-in direction.

The \emph{attraction}, the binding of bodies by radiation, is likewise emergent. It is the radiation-pressure shadow cast by a vacuum-excluding body in the active, growing mesh. It is a consequence of the knit's create move together with the open growing wake. It is not a base force.

\begin{principle}[The honesty rule]
If a phenomenon does not emerge from the eight atoms, report the negative. Do not bolt on a mechanism.
\end{principle}

\subsection{The tone is conserved}

This is the single most load-bearing fact in the base. Sum the tone of every site in the mesh. Treat each ternary value as a charge. The knit never changes that sum.

The collision only ever rearranges and balances tone. It never mints net charge from nothing. So the total tone is an exact conserved quantity, a constant of the whole crystal, fixed for all beats.

\begin{theorem}[The knit conserves total charge]
Let the total charge be the signed sum of the tones over every site of every dock in the mesh. The knit, both its collide phase and its stream phase, leaves this total charge unchanged at every beat. The total charge is therefore a constant of the motion. This was checked directly to machine precision on the $D_4$ sites.
\end{theorem}

What is conserved is the total, the signed sum. What is \emph{not} conserved is the arrangement, the order and pattern of where the charge sits, which is free to form, to flow, and to decay. The crystal keeps its books exactly while letting its patterns live and die.

This one fact does enormous downstream work. It is why the create move can only ever bring out balanced pairs, a $+1$ together with a $-1$, summing to zero, never a lone charge from nowhere. It forces newborn frontier docks to be born at peace, since a fresh dock with zero net charge is a dock at peace. And it gives the dilution toward peace as the dock count grows, because a fixed total charge spread over an ever-larger mesh thins out everywhere.

\subsection{The knit runs backward}

The collide phase is a reversible local map. Streaming is a bijection on docks. So the whole knit runs backward as exactly as it runs forward.

\begin{lemma}[Exact reversibility of the bulk]
One beat of the knit is an invertible permutation of states. Running forward $N$ beats and then running the exact inverse backward $N$ beats returns the start with zero error. This was verified directly.
\end{lemma}

Reversibility matters twice over. First, it is the structural reason the right quantities are conserved, the charge above among them. Second, and just as important, it means every state has a unique past. A reversible bulk has no forks looking backward, which is exactly what the cosmology of the model needs. Because the bulk is reversible, the arrow of time cannot come from the bulk. It must come from elsewhere, and the wake supplies it, as the later section on the arrow shows.

\subsection{What is base and what is emergent}

It helps to lay the line out plainly. The left column is the base, the right column is everything that appears only after coarse-graining.

\begin{table}[ht]
\centering
\begin{tabular}{@{}p{0.46\textwidth}p{0.46\textwidth}@{}}
\toprule
in the base (Layer 0) & not in the base (emergent, Layer 1 and above) \\
\midrule
ternary tone, $24$ directions, the collision table, the beat, growth & particles, mass, sound, light as fields \\
the finite symmetry $F_4$, of order $1152$ & continuous rotation, Lorentz symmetry, Lie groups \\
discrete charge conservation & smooth gravity, the $1/r$ force \\
the reversible bulk plus the open growing wake & the arrow, the felt direction of time, entropy growth \\
the active vacuum, the create move & the attraction, binding by radiation \\
the lean, the order $-1 < 0 < +1$ on the tone & the self, its identity, its agency, its bound body \\
\bottomrule
\end{tabular}
\caption{The base versus the emergent.}
\label{tab:base-vs-emergent}
\end{table}

Everything on the left could be stored exactly in a finite computer and stepped exactly. Everything on the right is a property of an average over many grains. The bridge between the two columns is the subject of the rest of the paper.

\subsection{The two-layer rule}

The base is discrete. The self and all recognizable physics are emergent. This single rule resolves the one apparent contradiction a careful reader meets early. The later soliton models use real numbers, so how can the theory claim to be built on a ternary tone and to be completely discrete? The answer is that there are two layers, and the real numbers live only in the upper one.

\paragraph{Layer 0, the base, fully discrete.} The state is tones, $\{-1, 0, +1\}$, on each site, and nothing else. Momentum is the sum over sites of the tone times its root vector, an integer vector, because tones and roots are both integers. The $24$ directions form a finite group, the unit Hurwitz quaternions, which is the binary tetrahedral group of $24$ elements. Composing two directions is quaternion multiplication, which stays inside the $24$, since the group is closed. This is finite-group arithmetic, as discrete as modular arithmetic on a finite set. At the base, to rotate by a fixed angle means to multiply by a fixed group element, a discrete operation on $24$ things, with no real angle anywhere. No real number ever appears at Layer 0.

\paragraph{Layer 1 and above, the self, emergent and effectively continuous.} Average many ternary sites over a region, and that average is a fine-grained quantity, effectively a real number. This is exactly how a magnet works. The smooth magnetization field is the average of discrete atomic spins. The atoms are discrete, the field is smooth. The self lives in this coarse-grained direction field, and the real numbers that show up in the soliton models are properties of that average, not of the base. A small real twist angle in an effective model is the coarse-grained stand-in for many discrete group-element twists averaged together. The decimal is an artifact of the averaging, not a base ingredient.

\begin{principle}[The reals are the song]
The reals are the song. The tones are the singers.
\end{principle}

The honest scorecard on the bridge from Layer 0 to Layer 1 is worth stating outright, since the value of the knit is that it lets us be exact about what has been shown.

\begin{table}[ht]
\centering
\begin{tabular}{@{}lp{0.5\textwidth}@{}}
\toprule
piece of the self & status \\
\midrule
soft radiation (sound) & demonstrated discretely, built and measured from ternary tones \\
bath damping (healing) & demonstrated discretely on the open substrate \\
identity (the winding) & demonstrated discretely at $24$-direction resolution \\
binding (the tight knot) & demonstrated discretely, the radiation-pressure shadow of a vacuum-excluding body, on integer tones and momentum, native on the $24$ sites \\
\bottomrule
\end{tabular}
\caption{The discrete demonstration of the self's parts.}
\label{tab:self-scorecard}
\end{table}

So all three faces of the self, agency, identity, and binding, are demonstrated on the discrete substrate. The binding is the radiation-pressure shadow of a vacuum-excluding body, fully discrete and reversible in the bulk. What remains is a decision and a long-time showcase, not a missing mechanism.

\subsection{The one seam the experiments cannot cross}

The base posits that the substrate \emph{feels}. That is the experiential posit, and it is the single premise the experiments do not test. They assume the substrate feels, and ask what follows.

\begin{result}[The model-territory seam]
Everything above the seam is structure that can be run and measured. Below the seam is the claim that the structure is made of feeling rather than merely describing it. The reversal, that experience is the ground and structure is what experience takes at scale, is a premise, not a result.
\end{result}

We mark the seam plainly so it cannot be smuggled past. What follows from the premise is a great deal, and all of it is checkable. The rest of the paper does the checking.
